\documentclass[CJK]{beamer}
\usepackage{CJKutf8}
\usepackage{beamerthemesplit}
\usetheme{Malmoe}
\useoutertheme[footline=authortitle]{miniframes}
\usepackage{amsmath}
\usepackage{amssymb}
\usepackage{graphicx}
\usepackage{color}
\usepackage{slashed}
\usepackage{simplewick}
\graphicspath{{../figures/}}
\def\be{\begin{equation}}
\def\ee{\nonumber\end{equation}}
\def\bea{\begin{eqnarray}}
\def\eea{\nonumber\end{eqnarray}}
\def\ii{{\dot{\imath}}}
\def\bch{\begin{CJK}{UTF8}{gbsn}}
\def\ech{\end{CJK}}
\def\bex{\begin{minipage}{0.3\textwidth}\includegraphics[width=1in]{jugelizi.png}\end{minipage}\begin{minipage}{0.6\textwidth}}
\def\eex{\end{minipage}}
\def\chtitle#1{\frametitle{\bch#1\ech}}
\def\skipline{{\vskip0.1in}}
\def\skiplines{{\vskip0.2in}}
\def\lagr{{\mathcal{L}}}
\def\hamil{{\mathcal{H}}}
\def\vecv{{\mathbf{v}}}
\def\vecx{{\mathbf{x}}}
\def\veck{{\mathbf{k}}}
\def\vecp{{\mathbf{p}}}
\def\vecn{{\mathbf{n}}}
\def\vecA{{\mathbf{A}}}
\def\vecP{{\mathbf{P}}}
\def\vecsigma{{\mathbf{\sigma}}}
\def\hatJn{{\hat{J_\vecn}}}
\def\hatJx{{\hat{J_x}}}
\def\hatJy{{\hat{J_y}}}
\def\hatJz{{\hat{J_z}}}
\def\hatj#1{\hat{J_{#1}}}
\def\hatphi{{\hat{\phi}}}
\def\hatq{{\hat{q}}}
\def\hatpi{{\hat{\pi}}}
\def\vel{\upsilon}
\def\Dint{{\mathcal{D}}}
\def\adag{{\hat{a}^\dagger}}
\def\bdag{{\hat{b}^\dagger}}
\def\cdag{{\hat{c}^\dagger}}
\def\ddag{{\hat{d}^\dagger}}
\def\hata{{\hat{a}}}
\def\hatb{{\hat{b}}}
\def\hatc{{\hat{c}}}
\def\hatd{{\hat{d}}}
\def\hatN{{\hat{N}}}
\def\hatH{{\hat{H}}}
\def\hatp{{\hat{p}}}
\def\Fup{{F^{\mu\nu}}}
\def\Fdown{{F_{\mu\nu}}}
\def\newl{\nonumber \\}
\def\SIkm{\mathrm{km}}
\def\SIyr{\mathrm{yr}}
\def\SIGyr{\mathrm{Gyr}}
\def\SIeV{\mathrm{eV}}
\def\SIGeV{\mathrm{GeV}}
\def\SIm{\mathrm{m}}
\def\SIcm{\mathrm{cm}}
\def\SIJ{\mathrm{J}}
\def\SIs{\mathrm{s}}
\def\SIkg{\mathrm{kg}}
\def\SIg{\mathrm{g}}
\def\vece{\mathrm{e}}
\def\bmat#1{\left(\begin{array}{#1}}
\def\emat{\end{array}\right)}
\def\bcase#1{\left\{\begin{array}{#1}}
\def\ecase{\end{array}\right.}
\def\calM{{\mathcal{M}}}
\def\calT{{\mathcal{T}}}
\def\calR{{\mathcal{R}}}
\def\barpsi{\bar{\psi}}
\def\baru{\bar{u}}
\def\barv{\bar{\upsilon}}
\def\bmini#1{\begin{minipage}{#1\textwidth}}
\def\emini{\end{minipage}}
\def\qeq{\stackrel{?}{=}}
\def\torder#1{\mathcal{T}\left(#1\right)}
\def\rorder#1{\mathcal{R}\left(#1\right)}


\title{Quantum Field Theory I \\ Lesson 06 - Angular Momentum}
\author{}
\date{}


\begin{document}

\begin{frame}
 
\begin{center}
\begin{Large}
\bch
量子场论 I 

{\vskip 0.3in}

第六课 角动量理论和矢量场的自旋

\ech
\end{Large}
\end{center}

\vskip 0.2in

\bch
课件下载
\ech
https://github.com/zqhuang/SYSU\_QFTI

\end{frame}



\begin{frame}
\chtitle{矩阵的函数简要回顾}
\bch
上节课末尾讲到，对任意的解析函数$f(x)$总能定义矩阵的函数$f(A)$。它有如下性质。
\begin{itemize}
\item{若$A$的本征值和本征矢为$(\lambda_1, \vecv_1)$; $(\lambda_2, \vecv_2)$;$\ldots$ 则$f(A)$有相同的本征矢，相应的本征值则变为$f(\lambda_1)$, $f(\lambda_2)$, $\ldots$}
\item{矩阵的函数满足复合函数规则，即：若解析函数$f$, $g$, $h$满足$h(x) = f(g(x))$，则对任意矩阵$A$有$h(A)=f(g(A))$。特别地，$(e^A)^n = e^{nA}$。}
\item{共轭转置和函数可以交换位置。即$\left(f(A)\right)^\dagger = f(A^\dagger)$。由此得出若$A$是厄米矩阵($A=A^\dagger$)，则$f(A)$也是厄米矩阵。注意厄米矩阵的本征值都是实数。}
\item{若两个矩阵$A$, $B$对易，则它们的任意函数$f(A)$与$g(B)$也对易。特别地，$ e^{A}e^{B} = e^{B}e^{A}=e^{A+B}$。}
\end{itemize}

\ech
\end{frame}

\begin{frame}
\chtitle{三维空间转动算符}
\bch
任意三维矢量$\vecv$绕方向$\vecn$（我们约定方向矢总是归一化的，即$\vecn^2=1$）旋转无穷小角度$\epsilon$总是可以写成：
$$\vecv \rightarrow \vecv - \ii\epsilon \hatJn \vecv = e^{-\ii\epsilon\hatJn}\vecv\,.$$
其中$\hatJn$是$3\times 3$的转动矩阵，它只跟$\vecn$有关。为了防止同三维空间指标$i$混淆，以后我们用符号$\ii$来表示虚数单位。

对于有限大小的转动角$\theta$，因为总是可以把转动划分成很多份的小转动，所以取足够大的$N$即有
$$\vecv\rightarrow (e^{-\ii(\theta/N)\hatJn})^N\vecv = e^{-\ii\theta \hatJn}\vecv$$
也就是说绕$\vecn$旋转$\theta$的操作仍然可以写成$e^{-\ii\theta \hatJn}$。算符$\hatJn$是对应于$\vecn$方向的转动算符，它不依赖于坐标系而存在。仅当我们需要写出$\hatJn$的具体矩阵表达式时，才需要选定一个坐标系。

\ech
\end{frame}

\begin{frame}
\chtitle{课堂互动}
\bch
对三个形成右手正交系的方向$\vecn_1$, $\vecn_2$, $\vecn_3$,对应的转动算符分别为$\hatj{1} $, $\hatj{2} $, $\hatj{3} $。试证明
\begin{itemize}
\item{对任意角度$\theta$，有$e^{-\ii\theta\hatj{3} }\vecn_1 =  \vecn_1 \cos\theta + \vecn_2 \sin\theta$, $e^{-\ii\theta\hatj{3} }\vecn_2 = \vecn_2 \cos\theta - \vecn_1 \sin\theta$, $e^{-\ii\theta\hatj{3}}\vecn_3 = \vecn_3$. }
\item{对任意角度$\theta$, 算符$e^{-\ii\theta\hatj{3} }$的三组本征矢量和本征值分别为$(\vecn_3, 1)$, $(\vecn_1+\ii\vecn_2, e^{-\ii\theta})$, $(\vecn_1-\ii\vecn_2, e^{\ii\theta})$。}
\item{根据上题结论，证明$\hatj{3}$的三组本征矢量和本征值分别$(\vecn_3, 0)$, $(\vecn_1+\ii\vecn_2, 1)$, $(\vecn_1-\ii\vecn_2, -1)$}.
\item{根据上题结论，证明$\hatj{3}\vecn_1=\ii\vecn_2$, $\hatj{3} \vecn_2 = -\ii\vecn_1$。通过置换$1,2,3$可以得到$\hatj{i}\vecn_j = \ii\epsilon_{ijk}\vecn_k$。其中$\epsilon_{ijk}$为完全反对称张量（当$ijk$为$123$的偶置换时为$\epsilon_{ijk}=1$，当$ijk$为$123$的奇置换时$\epsilon_{ijk}=-1$, $ijk$有重复指标时$\epsilon_{ijk} = 0$)。}
\item{$[\hatj{i}, \hatj{j} ] = \ii \epsilon_{ijk}\hatj{k} $}
\end{itemize}
\ech
\end{frame}



\begin{frame}
\chtitle{空间转动对称性和角动量守恒}
\bch
由于作用量是拉氏密度在全空间的积分，所以对固定的$\vecn$，整体空间的转动$e^{i\epsilon\hatJn}$保持作用量不变。对应的Noether定理中的守恒量就是角动量。

\skipline
首先，我们形式上把三维转动矩阵$\hatJn$拓写成四维张量$J^{\mu}_{\;\nu}$，规定$J^0_{\;\nu} = J^\mu_{\;0} = 0$，对$i$, $j$均不为零的情况$J^i_{\;j}$则是$\hatJn$矩阵的第$i$行第$j$列元素。(显然，这样定义的$J^\mu_{\;\nu}$并不是任意坐标变换下的张量，我们这里仅仅讨论三维空间的转动则没有问题。)由转动带来的坐标变动为
$$ \delta x^\mu =   \ii \epsilon J^\mu_{\;\nu} x^\nu$$

这样，对拉氏密度中的所有标量自由度$\phi$而言，$\phi$在转动后的$x^\mu$的位置的值，等于$\phi$在转动之前的坐标$x^\mu - \ii\epsilon J^\mu_{\;\nu} x^\nu$处的值。所以
$$\delta \phi = (-\ii \epsilon J^\mu_{\;\nu} x^\nu)\partial_\mu \phi$$

\ech
\end{frame}

\begin{frame}
\chtitle{空间转动对称性和角动量守恒}
\bch
于是我们得到
$$\frac{\delta \phi}{\delta\epsilon} = -\ii  J^\mu_{\;\nu} x^\nu \partial_\mu \phi$$.

\skipline
对矢量分量$A^\sigma$，则由于$A$本身也会发生旋转，所以多了额外一项。
$$\frac{\delta A^\sigma}{\delta\epsilon} = -\ii  J^\mu_{\;\nu} x^\nu\partial_\mu A^\sigma + \ii J^\sigma_{\;\nu}A^\nu$$.

\ech
\end{frame}

\begin{frame}
\chtitle{空间转动对称性和角动量守恒}
\bch
最后，拉氏密度作为一个标量，它的变化仅仅由坐标变动而引起：
$$ \delta \lagr = (-\ii \epsilon J^\mu_{\;\nu} x^\nu)\partial_\mu \lagr $$
另外，有
$$\partial_\mu(J^\mu_{\;\nu} x^\nu) = J^\mu_{\;\nu} \delta^\nu_\mu = J^\mu_\mu = 0 $$
最后一个等号是因为$\hatJn$的迹(Trace)是所有本征值的和($=1+(-1)+0 = 0$).

所以我们得到
$$\delta \lagr = \epsilon \partial_\mu (-\ii  J^\mu_{\;\nu} x^\nu \lagr) $$

也就是Noether定理中的$F^\mu = -\ii  J^\mu_{\;\nu} x^\nu \lagr$. 显然$F^0 = 0$.
\ech
\end{frame}

\begin{frame}
\chtitle{空间转动对称性和角动量守恒}
\bch
综上所述，我们得到Noether定理守恒荷密度：
$$j^0 =   \sum_q(\ii  J^\mu_{\;\nu} x^\nu \partial_\mu q) \frac{\partial \lagr}{\partial(\partial_0 q)} - \ii J^\sigma_{\;\nu}A^\nu \frac{\partial \lagr}{\partial(\partial_0 A^\sigma)}$$
其中 $q$取遍标量自由度和矢量所有分量自由度。因为$J^0_{\;\nu} = J^\mu_{\;0} = 0$，上式也可写成仅对空间坐标求和
$$j^0 =  \ii  J^i_{\;j} x^j \sum_q(\partial_i q )\frac{\partial \lagr}{\partial(\partial_0 q)} - \ii J^i_{\;j}A^j \frac{\partial \lagr}{\partial(\partial_0 A^i)}$$
\ech
\end{frame}

\begin{frame}
\chtitle{空间转动对称性和角动量守恒}
\bch
最后，利用上节课的结论$\frac{\lagr}{\partial(\partial_0 A^i)} = - F^0_{\; i}$，以及以前讨论过的空间平移带来的动量密度
$p_i = \sum_q(\partial_i q )\frac{\partial \lagr}{\partial(\partial_0 q)} $，有
$$j^0 =  \ii  J^i_{\;j} (x^j p_i + A^j F^0_{\;i} )$$
注意到$F^0_{\;i}$即为“电场强度”$\mathbf{E}$, 可以不失一般性地取右手正交坐标系，并设$\vecn$沿$\vecn_3$轴方向，根据$\hatj{3}\vecn_1 = i \vecn_2$, $\hatj{3}\vecn_2 = - i\vecn_1$，$\hatj{3}\vecn_3 = 0$, 上面的式子可以写成
$$j^0 = (\vecx \times \vecp + \mathbf{E} \times \mathbf{A})\cdot \vecn_3$$
注意因为取了$+---$的度规，$p_i$和$p^i$相差一个负号。上式第一项为轨道角动量（包含了所有自由度），第二项是由于空间转动时矢量会随之变化而带来的“自旋”。
\ech
\end{frame}

\begin{frame}
\chtitle{$U(1)$规范场的量子化}
\bch
现在考虑一个自由$U(1)$规范场
$$\lagr = -\frac{1}{4}\Fup\Fdown$$
怎样把它量子化？

\skipline
显然，由于矢量场存在自旋，它不能看成许多不相关的谐振子的叠加。我们首先需要建立自旋为$1$的粒子的量子理论。

大致思路是, 仍然把时间维特殊化，取库仑规范排除冗余自由度，然后把$A^\mu$场进行傅立叶变换。由于拉氏密度仍然是二次的，所以傅立叶变换后每一个傅立叶波矢$\veck$对拉氏量的贡献都是独立的。我们就只要对固定$\veck$的拉氏量研究其量子化。

\ech
\end{frame}

\begin{frame}
\chtitle{$U(1)$规范场的量子化}
\bch
我们先猜测一个形式解，由于$A^\mu$是实数场，它必须是互为共轭的两组算符的和。跟标量场不同的是，现在对每一个$\veck$，粒子均可能有自旋$s$：
$$\hat{\vecA}(x) = \frac{1}{(2\pi)^{3/2}} \int \sqrt{\frac{d^3\veck}{2\omega}} \mathbf{e}_{\veck s}\left(\hata_{\veck s} e^{-ik_\mu x^\mu} + \adag_{\veck s}e^{ik_\mu x^\mu}\right) $$

对每个$\veck$，自旋均按$\hat{J_\veck}$的本征矢进行分解。我们前面讨论过，本征值为零的本征矢$\mathbf{e}_0$就是$\veck$方向的单位矢量，取了库仑规范$\nabla\cdot A=0$之后$\vecA$不能有这个方向的分量，因此上面求和只要对自旋为$s=\pm 1$求和。
\ech
\end{frame}

\begin{frame}
\chtitle{第二次作业已经可以下载}
\bch
第二次作业10月10日交。

\skipline
下节课发第一次作业并简要解答。如果有时间我们详细讨论$U(1)$规范场的量子化。
\ech
\end{frame}

\end{document}
